\documentclass[11pt,a4paper]{article}
% Wspólny szablon dla kart supermocy (jedna karta = jedna strona A4)
% Użycie w karcie: \documentclass[11pt,a4paper]{article} % Wspólny szablon dla kart supermocy (jedna karta = jedna strona A4)
% Użycie w karcie: \documentclass[11pt,a4paper]{article} % Wspólny szablon dla kart supermocy (jedna karta = jedna strona A4)
% Użycie w karcie: \documentclass[11pt,a4paper]{article} \input{../preambula} \begin{document}

% Silnik tectonic (XeTeX): polskie znaki przez fontspec, nie przez inputenc/T1
\usepackage{fontspec}
\setmainfont{Arial}
\usepackage[polish]{babel}
\usepackage[a4paper,margin=1.5cm]{geometry}
\usepackage{amsmath,amssymb}
\usepackage{tikz}
\usetikzlibrary{arrows.meta,positioning,calc,shapes.geometric}
\usepackage[most]{tcolorbox}
\usepackage{enumitem}
\usepackage{booktabs}
\usepackage{xcolor}

\pagestyle{empty}
\setlength{\parindent}{0pt}
\setlength{\parskip}{0.4em}

\definecolor{cel}{HTML}{1F4E79}
\definecolor{regula}{HTML}{E8F1FA}
\definecolor{pulapka}{HTML}{FDECEA}
\definecolor{okres}{HTML}{C0392B}
\definecolor{skonczony}{HTML}{1E8449}

% Ramka na regułę
\newtcolorbox{regula}[1][]{
  colback=regula, colframe=cel, boxrule=0.8pt, arc=2pt,
  left=6pt, right=6pt, top=4pt, bottom=4pt, title=#1,
  fonttitle=\bfseries, coltitle=white
}

% Ramka na pułapkę z kartkówki
\newtcolorbox{pulapka}[1][]{
  colback=pulapka, colframe=okres, boxrule=0.8pt, arc=2pt,
  left=6pt, right=6pt, top=4pt, bottom=4pt, title=#1,
  fonttitle=\bfseries, coltitle=white
}

% Nagłówek karty: numer, tytuł, cel
\newcommand{\naglowekKarty}[3]{%
  {\LARGE\bfseries\color{cel} Karta #1: #2}\\[2pt]
  {\large Cel: #3}\par
  \vspace{0.3em}\noindent\rule{\linewidth}{1pt}\par
}
 \begin{document}

% Silnik tectonic (XeTeX): polskie znaki przez fontspec, nie przez inputenc/T1
\usepackage{fontspec}
\setmainfont{Arial}
\usepackage[polish]{babel}
\usepackage[a4paper,margin=1.5cm]{geometry}
\usepackage{amsmath,amssymb}
\usepackage{tikz}
\usetikzlibrary{arrows.meta,positioning,calc,shapes.geometric}
\usepackage[most]{tcolorbox}
\usepackage{enumitem}
\usepackage{booktabs}
\usepackage{xcolor}

\pagestyle{empty}
\setlength{\parindent}{0pt}
\setlength{\parskip}{0.4em}

\definecolor{cel}{HTML}{1F4E79}
\definecolor{regula}{HTML}{E8F1FA}
\definecolor{pulapka}{HTML}{FDECEA}
\definecolor{okres}{HTML}{C0392B}
\definecolor{skonczony}{HTML}{1E8449}

% Ramka na regułę
\newtcolorbox{regula}[1][]{
  colback=regula, colframe=cel, boxrule=0.8pt, arc=2pt,
  left=6pt, right=6pt, top=4pt, bottom=4pt, title=#1,
  fonttitle=\bfseries, coltitle=white
}

% Ramka na pułapkę z kartkówki
\newtcolorbox{pulapka}[1][]{
  colback=pulapka, colframe=okres, boxrule=0.8pt, arc=2pt,
  left=6pt, right=6pt, top=4pt, bottom=4pt, title=#1,
  fonttitle=\bfseries, coltitle=white
}

% Nagłówek karty: numer, tytuł, cel
\newcommand{\naglowekKarty}[3]{%
  {\LARGE\bfseries\color{cel} Karta #1: #2}\\[2pt]
  {\large Cel: #3}\par
  \vspace{0.3em}\noindent\rule{\linewidth}{1pt}\par
}
 \begin{document}

% Silnik tectonic (XeTeX): polskie znaki przez fontspec, nie przez inputenc/T1
\usepackage{fontspec}
\setmainfont{Arial}
\usepackage[polish]{babel}
\usepackage[a4paper,margin=1.5cm]{geometry}
\usepackage{amsmath,amssymb}
\usepackage{tikz}
\usetikzlibrary{arrows.meta,positioning,calc,shapes.geometric}
\usepackage[most]{tcolorbox}
\usepackage{enumitem}
\usepackage{booktabs}
\usepackage{xcolor}

\pagestyle{empty}
\setlength{\parindent}{0pt}
\setlength{\parskip}{0.4em}

\definecolor{cel}{HTML}{1F4E79}
\definecolor{regula}{HTML}{E8F1FA}
\definecolor{pulapka}{HTML}{FDECEA}
\definecolor{okres}{HTML}{C0392B}
\definecolor{skonczony}{HTML}{1E8449}

% Ramka na regułę
\newtcolorbox{regula}[1][]{
  colback=regula, colframe=cel, boxrule=0.8pt, arc=2pt,
  left=6pt, right=6pt, top=4pt, bottom=4pt, title=#1,
  fonttitle=\bfseries, coltitle=white
}

% Ramka na pułapkę z kartkówki
\newtcolorbox{pulapka}[1][]{
  colback=pulapka, colframe=okres, boxrule=0.8pt, arc=2pt,
  left=6pt, right=6pt, top=4pt, bottom=4pt, title=#1,
  fonttitle=\bfseries, coltitle=white
}

% Nagłówek karty: numer, tytuł, cel
\newcommand{\naglowekKarty}[3]{%
  {\LARGE\bfseries\color{cel} Karta #1: #2}\\[2pt]
  {\large Cel: #3}\par
  \vspace{0.3em}\noindent\rule{\linewidth}{1pt}\par
}


\begin{document}

\naglowekKarty{4}{Porównywanie ułamków}{wpisać $>$, $<$ lub $=$ między dwiema liczbami wymiernymi}

\vspace{0.3em}
\begin{regula}[REGUŁA – WYBIERZ JEDNĄ Z CZTERECH DRÓG]
\begin{enumerate}[leftmargin=*, itemsep=2pt]
  \item \textbf{Dziesiętne:} wyrównaj liczbę cyfr po przecinku zerami. $0{,}6$ vs $0{,}57$: $0{,}60 > 0{,}57$.
  \item \textbf{Okres:} rozpisz okres, aż będzie widać różnicę. $0{,}(4)=0{,}444\ldots > 0{,}400 = 0{,}4$.
  \item \textbf{Motylek} (ułamki zwykłe): pomnóż na krzyż. $\dfrac{a}{b}$ vs $\dfrac{c}{d}$ porównujemy $a\cdot d$ z $c\cdot b$.
  \item \textbf{Ta sama góra / ten sam dół:} przy tym samym liczniku większy jest mniejszy mianownik ($\dfrac{3}{7}>\dfrac{3}{8}$).
\end{enumerate}
\end{regula}

\vspace{0.4em}
\textbf{Dlaczego motylek działa? Druga droga: oś liczbowa.} Porównanie to pytanie, która liczba leży bardziej na prawo. Przybliżmy fragment osi między $0{,}6$ a $0{,}8$:

\begin{center}
\begin{tikzpicture}[>=Stealth, x=2.2cm]
  % Oś od 0,6 do 0,8; 1 jednostka na osi = 0,2 wartości (x=24cm, więc 0,6 -> 0, 0,8 -> 4.8cm)
  \pgfmathsetmacro{\sc}{20}
  \pgfmathsetmacro{\zero}{0.6}
  \draw[->] (-0.2,0) -- (4.6,0) node[right, font=\small] {$x$};
  \foreach \v in {0.6,0.65,0.7,0.75,0.8} {
    \pgfmathsetmacro{\px}{(\v-0.6)*\sc}
    \draw (\px,0.1) -- (\px,-0.1);
    \node[below, font=\small] at (\px,-0.1) {\pgfmathprintnumber[fixed,precision=2]{\v}};
  }
  % 7/10 = 0,70 oraz 5/7 ~ 0,714
  \pgfmathsetmacro{\pA}{(0.7-0.6)*\sc}
  \pgfmathsetmacro{\pB}{(5/7-0.6)*\sc}
  \fill[cel] (\pA,0) circle (2.2pt) node[above left=2pt, font=\small\bfseries, color=cel] {$\frac{7}{10}=0{,}7$};
  \fill[okres] (\pB,0) circle (2.2pt) node[above right=2pt, font=\small\bfseries, color=okres] {$\frac{5}{7}\approx 0{,}714$};
  \draw[dashed, okres] (\pB,0) -- (\pB,-0.5);
  \node[font=\small, color=okres] at (\pB,-0.7) {$\frac{5}{7}$ leży bardziej na prawo};
  \node[font=\small, anchor=west] at (-0.2,1.4) {Motylek: $5\cdot 10=50$, $7\cdot 7=49$, więc $50>49$, czyli $\dfrac{5}{7}>\dfrac{7}{10}$};
\end{tikzpicture}
\end{center}

\vspace{0.2em}
\textbf{Liczby ujemne (termometr):} im bliżej zera, tym większa liczba. $-0{,}2$ jest bliżej zera niż $-0{,}25$, więc $-0{,}2>-0{,}25$.

\vspace{0.2em}
\textbf{Przykłady rozstrzygnięte:}
\begin{center}
\renewcommand{\arraystretch}{1.4}
\begin{tabular}{@{}lll@{}}
\toprule
Porównanie & Metoda & Wynik \\
\midrule
$0{,}6$ vs $0{,}57$ & wyrównanie zerami: $0{,}60$ vs $0{,}57$ & $0{,}6>0{,}57$ \\
$0{,}(4)$ vs $0{,}4$ & rozpisz okres: $0{,}444\ldots$ vs $0{,}400$ & $0{,}(4)>0{,}4$ \\
$\dfrac{5}{7}$ vs $\dfrac{7}{10}$ & motylek: $50$ vs $49$ & $\dfrac{5}{7}>\dfrac{7}{10}$ \\
$\dfrac{3}{7}$ vs $\dfrac{3}{8}$ & ten sam licznik & $\dfrac{3}{7}>\dfrac{3}{8}$ \\
$-0{,}2$ vs $-0{,}25$ & termometr: bliżej zera & $-0{,}2>-0{,}25$ \\
\bottomrule
\end{tabular}
\end{center}

\begin{pulapka}[PUŁAPKI Z KARTKÓWKI]
\begin{enumerate}[leftmargin=*, itemsep=2pt]
  \item $0{,}4\neq 0{,}(4)$. Okres to nieskończony ciąg czwórek, więc $0{,}(4)$ jest większe.
  \item Więcej cyfr nie znaczy większa liczba: $0{,}7>0{,}68$, bo po wyrównaniu zerami $0{,}70>0{,}68$.
  \item Przy liczbach ujemnych odwróć intuicję: większa wartość bezwzględna oznacza \emph{mniejszą} liczbę.
\end{enumerate}
\end{pulapka}

\end{document}
